% !PRINTED 513
\weberpage{513}

If in (2) one replaces $v$ by another class invariant of the same discriminant $4D$ and of the same genus, then $u$ too passes over into another class invariant of the discriminant $D$, but likewise remains in the same genus.

Accordingly, if in (2) one lets $v$ run through the class invariants of \gesperrt{one genus of the discriminant} $4D$, then $u$ runs through the invariants of the corresponding genus of the discriminant $D$ [taking each of these values three times, and for $D\equiv 1\pmod 8$ only once], and the symmetric functions of these $u$ are contained in the same domain of rationality as the quantities $v$. If therefore the class function $H_{-4D}(v)$ is decomposable into factors according to the genera, the same holds of $H_{-D}(u)^{1)}$.

\textbf{§ 138. Splitting of the Class Equation according to the Genera.}

In order to apply Theorem 1., § 136, one sets
\weq{1}{D=B^2-4AC=-4m,}
and decomposes $m$ into two factors
\weq{2}{m=m'm'',}
\gesperrt{where it is assumed that $m''$ is odd and without square divisor.}

Now in Theorem 1., § 136, we make the assumption:
\weq{3}{\begin{aligned}
x&=4,\quad y=2(4m'-m''),\\
n&=(4m'+m'')^2.
\end{aligned}}
\weq{4}{\begin{aligned}
\alpha&=2B+4m'-m'',\\
\alpha m''&=4AC-(B-m'')^2,\\
\beta&=4A.
\end{aligned}}
\weq{5}{\begin{aligned}
\alpha+\beta\omega&=2i\sqrt{m'm''}+4m'-m'',\\
&=(2\sqrt{m'}+i\sqrt{m''})^2,
\end{aligned}}
% SOURCE-ORIGINAL-SIC: Equation (5) prints 2i sqrt(m'm'') in its first line, while the square in its second line expands with 4i sqrt(m'm''); the printed coefficient 2 is preserved.

\footnotetext[1]{The carrying out of the corresponding considerations for an odd $D$ would indeed also be possible, but would make numerous distinctions and prolixities necessary. Here is a point where the \gesperrt{Gauss}ian notation and distinction of forms of the first and second kind, of which I still made use in the first edition, would bring with it a certain simplification of the expression. This is connected with the exceptional position which the number 2 occupies in the whole theory of elliptic functions, and which in the \gesperrt{Weierstrass}ian theory recedes somewhat, but yet does not entirely disappear. On the other hand, this singling out of the number 2 is again the source of great simplifications, especially in the numerical results.}
%CONT>>
